%!TeX program = xelatex
\documentclass{SYSUReport}
\usepackage{tabularx} % 在导言区添加此行
\usepackage{float}
\usepackage{amsmath}
% Listings configuration for CUDA C++
\lstdefinestyle{cuda_cpp}{
    language=C++,
    basicstyle=\ttfamily\footnotesize,
    keywordstyle=\color{blue},
    commentstyle=\color{green!40!black},
    stringstyle=\color{purple},
    numbers=left,
    numberstyle=\tiny\color{gray},
    breaklines=true,
    breakatwhitespace=true,
    tabsize=2,
    captionpos=b, % Caption at the bottom
    escapeinside={\%*}{*)}, % Allows LaTeX commands inside listings
    morekeywords={__global__, __shared__, threadIdx, blockIdx, blockDim, gridDim, cudaMalloc, cudaMemcpy, cudaFree, __syncthreads, dim3}, % Add CUDA specific keywords
    backgroundcolor=\color{gray!10},
    frame=single,
    framerule=0.5pt,
    rulecolor=\color{gray!40}
}
% 根据个人情况修改
\headl{}
\headc{}
\headr{并行程序设计与算法实验}
\lessonTitle{并行程序设计与算法实验}
\reportTitle{Lab10-CUDA并行矩阵乘法}
\stuname{林宏宇}
\stuid{23320093}
\inst{计算机学院}
\major{计算机科学与技术}
\date{2025年6月13日}

\begin{document}

% =============================================
% Part 1: 封面
% =============================================
\cover
\thispagestyle{empty} % 首页不显示页码
\clearpage

% % =============================================
% % Part 4： 正文内容
% % =============================================
% % 重置页码，并使用阿拉伯数字
% \pagenumbering{arabic}
% \setcounter{page}{1}

%%可选择这里也放一个标题
%\begin{center}
%    \title{ \Huge \textbf{{标题}}}
%\end{center}

\section{实验目的}
\begin{itemize}
    \item 理解CUDA编程模型(Grid、Block、Thread)及其在矩阵乘法中的应用。
    \item 学习GPU内存优化技术。
\end{itemize}

\section{实验内容}
\begin{itemize}
\item 实现基础矩阵乘法
\item 优化矩阵乘法：共享内存，分块技术
\item 测量不同实现的运行时间
\end{itemize}

\section{实验原理}

\subsection{CUDA编程模型与矩阵乘法}

CUDA编程模型将计算任务组织为Grid-Block-Thread三级层次结构。在矩阵乘法$C = A \times B$（其中$A$为$M\times K$矩阵，$B$为$K\times N$矩阵，$C$为$M\times N$矩阵）中，每个输出元素$C_{ij} = \sum_{k=0}^{K-1} A_{ik} \cdot B_{kj}$的计算是独立的，天然适合并行化。

实验采用三种不同粒度的优化策略：
\begin{enumerate}
    \item \textbf{朴素实现（Naive）}：每个线程直接从全局内存读取数据，计算一个输出元素。全局内存访问延迟通常为数百个时钟周期，且存在大量重复读取——$A$的同一行被$N$个线程重复读取，$B$的同一列被$M$个线程重复读取。
    \item \textbf{共享内存分块（Shared Memory Tiling）}：将矩阵划分为子块（tile），每次将一个tile的数据从全局内存批量加载到共享内存（on-chip，延迟约为20--30个时钟周期），块内所有线程可重复使用这些数据，实现"一次加载，多次计算"。
    \item \textbf{寄存器分块（Register Tiling）}：在共享内存分块基础上，每个线程进一步使用寄存器缓存一个小型子块（如$4\times 4$），在寄存器中完成多次乘累加后再写回，将计算访存比从0.5提升至2.0（以tile\_size=4为例）。
\end{enumerate}

\subsection{数据划分策略}

实验实现了三种数据/任务划分方式，分别对应不同的Grid和Block组织方式：

\begin{itemize}
    \item \textbf{按行划分（Row-wise）}：Grid为1维，每个thread block负责一组完整的行（所有$N$列）。Block内线程沿列方向并行，沿$K$方向分块迭代。此方式对$A$矩阵访问具有良好的空间局部性（行优先存储），但对$B$矩阵的列访问需要跨步读取。
    \item \textbf{按列划分（Column-wise）}：Grid为1维，每个thread block负责一组完整的列（所有$M$行）。此方式对$B$矩阵访问有良好的空间局部性，但$A$矩阵需要按行读取（恰好与行优先存储匹配），而$B$的列访问存在严重的非合并访存问题。
    \item \textbf{按数据块划分（Block-wise / 2D Tiling）}：Grid为2维，每个thread block负责$C$矩阵的一个2维子块。这是最标准且最高效的CUDA矩阵乘法实现方式，兼顾了$A$和$B$的访存合并性，同时提供了两个维度的并行度。
\end{itemize}

\section{实验实现}

\begin{itemize}
    \item 实现基础CUDA矩阵乘法（朴素全局内存版本）
    \item 优化矩阵乘法：基于共享内存的分块优化
    \item 进一步优化：基于寄存器的分块优化（寄存器复用）
    \item 实现三种数据划分策略：按行划分、按列划分、按数据块划分
    \item 测试不同矩阵规模（512、1024、2048）和线程块大小（8$\times$8、16$\times$16、32$\times$32）
    \item 测量并对比不同实现的运行时间，分析性能差异原因
\end{itemize}

\subsection{朴素实现——按数据块划分}

\begin{lstlisting}[style=cuda_cpp, caption={朴素矩阵乘法（全局内存，2D Grid + 2D Block）}, label={lst:naive}]
__global__ void matmul_naive_block(const float *__restrict__ A,
                                    const float *__restrict__ B,
                                    float *__restrict__ C,
                                    int M, int N, int K) {
    int row = blockIdx.y * blockDim.y + threadIdx.y;
    int col = blockIdx.x * blockDim.x + threadIdx.x;

    if (row < M && col < N) {
        float sum = 0.0f;
        for (int k = 0; k < K; k++) {
            sum += A[row * K + k] * B[k * N + col];
        }
        C[row * N + col] = sum;
    }
}
\end{lstlisting}

朴素实现中，每个线程计算一个输出元素，遍历整个$K$维度。主要问题是：（1）每个线程都需要从全局内存读取$2K$个浮点数，访存量极大；（2）$A$的同一行被$N$个线程重复读取，$B$的同一列被$M$个线程重复读取，全局内存带宽被严重浪费。

\subsection{共享内存分块实现}

\begin{lstlisting}[style=cuda_cpp, caption={共享内存分块矩阵乘法}, label={lst:shared}]
template <int BLOCK_SIZE>
__global__ void matmul_shared_block(const float *__restrict__ A,
                                     const float *__restrict__ B,
                                     float *__restrict__ C,
                                     int M, int N, int K) {
    __shared__ float As[BLOCK_SIZE][BLOCK_SIZE];
    __shared__ float Bs[BLOCK_SIZE][BLOCK_SIZE];

    int row = blockIdx.y * BLOCK_SIZE + threadIdx.y;
    int col = blockIdx.x * BLOCK_SIZE + threadIdx.x;

    float sum = 0.0f;
    for (int tk = 0; tk < K; tk += BLOCK_SIZE) {
        // 协作加载A的子块到共享内存
        if (row < M && (tk + threadIdx.x) < K)
            As[threadIdx.y][threadIdx.x] = A[row * K + tk + threadIdx.x];
        else
            As[threadIdx.y][threadIdx.x] = 0.0f;

        // 协作加载B的子块到共享内存
        if ((tk + threadIdx.y) < K && col < N)
            Bs[threadIdx.y][threadIdx.x] = B[(tk + threadIdx.y) * N + col];
        else
            Bs[threadIdx.y][threadIdx.x] = 0.0f;

        __syncthreads();  // 保证整个block的数据加载完成

        // 利用共享内存进行乘累加
        for (int k = 0; k < BLOCK_SIZE; k++) {
            sum += As[threadIdx.y][k] * Bs[k][threadIdx.x];
        }
        __syncthreads();  // 保证下一轮加载前所有线程用完当前数据
    }

    if (row < M && col < N)
        C[row * N + col] = sum;
}
\end{lstlisting}

共享内存分块的核心思想：将$K$维度切分为多个大小为\texttt{BLOCK\_SIZE}的tile。每个thread block协作将$A$和$B$的一个tile加载到共享内存中，然后所有线程使用共享内存中的数据进行计算。数据重用率显著提高——每个从全局内存加载的元素被block内\texttt{BLOCK\_SIZE}个线程共享使用。

\subsection{寄存器分块实现}

\begin{lstlisting}[style=cuda_cpp, caption={寄存器分块矩阵乘法（关键计算部分）}, label={lst:reg}]
// TILE_DIM = 32 / BLOCK_SIZE, 即每个线程计算 TILE_DIM×TILE_DIM 个元素
__shared__ float As[BLOCK_SIZE * TILE_DIM][BLOCK_SIZE];
__shared__ float Bs[BLOCK_SIZE][BLOCK_SIZE * TILE_DIM];

// 寄存器中的累加器
float c[TILE_DIM][TILE_DIM] = {{0.0f}};
float a_reg[TILE_DIM], b_reg[TILE_DIM];

for (int tk = 0; tk < K; tk += BLOCK_SIZE) {
    // 协作加载：每个线程加载TILE_DIM个元素到更大的共享内存tile
    for (int i = 0; i < TILE_DIM; i++) {
        As[threadIdx.y * TILE_DIM + i][threadIdx.x] = ...;
        Bs[threadIdx.y][threadIdx.x * TILE_DIM + i] = ...;
    }
    __syncthreads();

    // 利用寄存器进行高密度计算
    for (int k = 0; k < BLOCK_SIZE; k++) {
        for (int i = 0; i < TILE_DIM; i++)
            a_reg[i] = As[threadIdx.y * TILE_DIM + i][k];
        for (int j = 0; j < TILE_DIM; j++)
            b_reg[j] = Bs[k][threadIdx.x * TILE_DIM + j];
        // 外积累加: TILE_DIM^2次乘加, 仅2*TILE_DIM次寄存器加载
        for (int i = 0; i < TILE_DIM; i++)
            for (int j = 0; j < TILE_DIM; j++)
                c[i][j] += a_reg[i] * b_reg[j];
    }
    __syncthreads();
}
\end{lstlisting}

寄存器分块的关键优化：每个线程处理的输出元素从1个增加到\texttt{TILE\_DIM}$\times$\texttt{TILE\_DIM}个。在每次从共享内存加载$2\times\texttt{TILE\_DIM}$个值到寄存器后，可完成$\texttt{TILE\_DIM}^2$次乘加运算。以\texttt{TILE\_DIM}=4为例，计算访存比从0.5（朴素/共享内存）提升至$16/(4+4)=2.0$，提升4倍。设计中保持\texttt{BLOCK\_SIZE}$\times$\texttt{TILE\_DIM}=32不变，使共享内存tile大小恒定为$32\times 32$（4KB），确保在不同block大小下都有合理的共享内存使用量。

\section{实验结果与分析}

\subsection{性能测试数据}

实验在NVIDIA Tesla V100S GPU上进行，矩阵规模$M=N=K \in \{512, 1024, 2048\}$，线程块大小为$8\times 8$、$16\times 16$和$32\times 32$。实验具体流程如下：（1）使用cuRAND库在GPU端生成服从均匀分布$U(0,1)$的随机矩阵$A$和$B$，固定随机种子以保证可复现性；（2）分别调用9种不同kernel（3种实现$\times$3种划分策略）执行矩阵乘法，每种配置先执行1次warm-up以消除冷启动偏差，再执行5次正式运行取平均时间；（3）将GPU计算结果传回主机，与CPU端串行矩阵乘法结果逐元素比较，验证条件为相对误差$< 10^{-3}$或绝对误差$< 10^{-3}$（对接近零值的情形）。所有81组测试配置均通过正确性验证。运行时间单位为毫秒（ms）。

\begin{table}[H]
\centering
\caption{按行划分：性能对比 (时间单位: ms)}
\label{tab:perf_row}
\begin{tabular}{|c|c|c|c|c|}
\hline
矩阵规模 (N) & 线程块大小 & 朴素实现 & 基于共享内存优化 & 基于寄存器分块优化 \\
\hline
\multirow{3}{*}{512} & 8$\times$8 & 1.292 & 0.728 & 1.326 \\
\cline{2-5}
& 16$\times$16 & 0.667 & 0.395 & 0.435 \\
\cline{2-5}
& 32$\times$32 & 0.534 & 0.352 & 0.359 \\
\hline
\multirow{3}{*}{1024} & 8$\times$8 & 7.659 & 5.995 & 14.028 \\
\cline{2-5}
& 16$\times$16 & 10.189 & 4.754 & 4.935 \\
\cline{2-5}
& 32$\times$32 & 5.151 & 1.669 & 1.640 \\
\hline
\multirow{3}{*}{2048} & 8$\times$8 & 45.625 & 20.068 & 33.432 \\
\cline{2-5}
& 16$\times$16 & 21.825 & 10.679 & 11.502 \\
\cline{2-5}
& 32$\times$32 & 11.768 & 6.717 & 6.803 \\
\hline
\end{tabular}
\end{table}

\begin{table}[H]
\centering
\caption{按列划分：性能对比 (时间单位: ms)}
\label{tab:perf_col}
\begin{tabular}{|c|c|c|c|c|}
\hline
矩阵规模 (N) & 线程块大小 & 朴素实现 & 基于共享内存优化 & 基于寄存器分块优化 \\
\hline
\multirow{3}{*}{512} & 8$\times$8 & 0.867 & 0.767 & 1.489 \\
\cline{2-5}
& 16$\times$16 & 1.422 & 0.813 & 1.097 \\
\cline{2-5}
& 32$\times$32 & 5.466 & 2.211 & 2.210 \\
\hline
\multirow{3}{*}{1024} & 8$\times$8 & 10.765 & 6.283 & 13.339 \\
\cline{2-5}
& 16$\times$16 & 11.288 & 6.271 & 9.348 \\
\cline{2-5}
& 32$\times$32 & 46.211 & 16.714 & 16.717 \\
\hline
\multirow{3}{*}{2048} & 8$\times$8 & 39.278 & 22.740 & 36.521 \\
\cline{2-5}
& 16$\times$16 & 45.516 & 23.227 & 21.404 \\
\cline{2-5}
& 32$\times$32 & 88.361 & 35.918 & 35.919 \\
\hline
\end{tabular}
\end{table}

\begin{table}[H]
\centering
\caption{按数据块划分：性能对比 (时间单位: ms)}
\label{tab:perf_block}
\begin{tabular}{|c|c|c|c|c|}
\hline
矩阵规模 (N) & 线程块大小 & 朴素实现 & 基于共享内存优化 & 基于寄存器分块优化 \\
\hline
\multirow{3}{*}{512} & 8$\times$8 & 0.189 & 0.123 & \textbf{0.053} \\
\cline{2-5}
& 16$\times$16 & 0.129 & 0.080 & \textbf{0.052} \\
\cline{2-5}
& 32$\times$32 & 0.140 & 0.083 & 0.081 \\
\hline
\multirow{3}{*}{1024} & 8$\times$8 & 1.342 & 0.910 & \textbf{0.303} \\
\cline{2-5}
& 16$\times$16 & 0.943 & 0.584 & \textbf{0.307} \\
\cline{2-5}
& 32$\times$32 & 0.891 & 0.549 & 0.523 \\
\hline
\multirow{3}{*}{2048} & 8$\times$8 & 19.286 & 8.155 & \textbf{2.893} \\
\cline{2-5}
& 16$\times$16 & 8.134 & 4.804 & \textbf{2.457} \\
\cline{2-5}
& 32$\times$32 & 7.746 & 4.613 & 4.140 \\
\hline
\end{tabular}
\end{table}

\subsection{分析性能差异的原因}

\begin{itemize}
    \item \textbf{结合CUDA内存模型和矩阵乘法原理，分析造成观察到的性能差异的可能原因。}\par
    回答：性能差异的根本原因在于GPU内存层次结构的带宽和延迟差异：
    \begin{enumerate}
        \item \textbf{全局内存带宽瓶颈}：V100的全局内存（HBM2）带宽约900GB/s，但延迟高达数百周期。朴素实现中每个输出元素的每次乘加都需要2次全局内存读取（共$2K$次），严重受限于内存带宽。由表\ref{tab:perf_block}可见，N=2048时朴素实现（block=16$\times$16）耗时8.134ms，而共享内存版本仅需4.804ms，加速约1.69$\times$；N=1024时朴素实现耗时0.943ms，共享内存版本为0.584ms，加速约1.61$\times$。这是因为朴素实现需读取约$2N^3$个浮点数，例如N=2048时约$17.2\times 10^9$次全局内存读取（约68.7GB），严重超出L2缓存容量，实际有效带宽远低于峰值。
        \item \textbf{共享内存减少全局访存}：共享内存带宽约13.8TB/s（V100），延迟约20--30周期。通过分块将数据预加载到共享内存，每个block内的\texttt{BLOCK\_SIZE}个线程共享同一份数据，全局访存减少约\texttt{BLOCK\_SIZE}倍。从表\ref{tab:perf_block}可以看出共享内存优化的显著效果：N=2048、block=32$\times$32时，朴素实现7.746ms降至共享内存的4.613ms（加速1.68$\times$）；N=1024时，相同配置下从0.891ms降至0.549ms（加速1.62$\times$）。同时，共享内存的bank结构支持高效的并行访问。
        \item \textbf{寄存器进一步提升计算密度}：寄存器延迟仅1个周期，是最快的存储层次。寄存器分块将多次共享内存读取合并为一次寄存器加载，增加每字节加载对应的计算量（计算访存比），使compute bound的潜力得以发挥。表\ref{tab:perf_block}的数据充分说明了这一点：N=2048、block=16$\times$16时，寄存器分块仅需2.457ms，相比共享内存的4.804ms加速1.96$\times$，相比朴素实现的8.134ms加速3.31$\times$。对于block=8$\times$8（TILE\_DIM=4），寄存器分块（2.893ms）相比共享内存（8.155ms）加速达2.82$\times$，体现了高计算访存比带来的显著收益。
        \item \textbf{访存合并性}：按列划分性能极差的原因是线程束（warp）内的32个线程访问的是步长很大的非连续地址（行优先存储下，同一列的元素在内存中相距$K$个元素），无法合并为一个内存事务，导致实际有效带宽远低于峰值。对比表\ref{tab:perf_row}和表\ref{tab:perf_col}可知：N=2048、block=32$\times$32时，按行划分共享内存版本耗时为6.717ms，而按列划分为35.918ms，性能相差5.35倍；N=2048、block=16$\times$16时差异更明显（10.679ms vs 23.227ms，相差2.17倍）。此外，按列划分中block=32$\times$32反而比16$\times$16更慢（朴素实现：88.361ms vs 45.516ms），这与常规预期相反，原因是大block加剧了非合并访存的延迟惩罚。
    \end{enumerate}

    \item \textbf{如何选择合适的线程块大小以提高占用率？}\par
    回答：选择线程块大小的基本原则如下：
    \begin{enumerate}
        \item \textbf{线程数应为warp大小（32）的整数倍}：确保warp充分利用，避免线程浪费。本实验使用的8$\times$8=64、16$\times$16=256、32$\times$32=1024均满足此要求。
        \item \textbf{考虑SM资源限制}：V100每个SM最多驻留2048个线程、64个warp、以及最多32个线程块。线程块过小（如8$\times$8=64线程）会导致SM上线程总数受限于block数量上限（32 blocks $\times$ 64 threads = 2048 threads，刚好达到上限），但此时每个线程的寄存器使用量是限制因素。线程块过大（如32$\times$32=1024线程）则每个SM只能驻留2个block（2048/1024=2），限制了延迟隐藏能力。
        \item \textbf{中等大小的线程块（128--256线程）通常最优}：在本实验中，16$\times$16线程块通常能够在占用率、共享内存利用率和访存效率之间取得较好的平衡，因此整体表现较为稳定。实验结果表明，16$\times$16线程块在大多数测试中具有较好的综合性能，但并非所有情况下都是最优。实际性能还受到寄存器使用量、共享内存占用率以及矩阵规模等因素影响。例如，N=2048、按数据块划分时，16$\times$16的寄存器分块（2.457ms）优于32$\times$32（4.140ms），但在共享内存版本中32$\times$32（4.613ms）略优于16$\times$16（4.804ms），说明不同优化策略下的最优block大小并不相同。
        \item \textbf{结合寄存器分块时需重新权衡}：寄存器分块增加了每线程的寄存器使用量，降低了SM上可驻留的block数。此时较小线程块（如8$\times$8）配合较大的TILE\_DIM（如4），可能获得更好的计算访存比，弥补线程数不足的劣势。
    \end{enumerate}

    \item \textbf{思考如果按不同的方式划分（例如，按行、列、数据块划分），可能会对性能和实现复杂度带来什么影响？}\par
    回答：

\begin{table}[H]
\centering
\caption{三种划分策略对比总结}
\label{tab:comparison}
\begin{tabularx}{\textwidth}{|l|X|X|X|}
\hline
特性 & 按行划分 & 按列划分 & 按数据块划分 \\
\hline
Grid维度 & 1D & 1D & 2D \\
\hline
并行度 & 有限（=行数/block\_size） & 有限（=列数/block\_size） & 充分（$\approx N^2$/tile\_size$^2$） \\
\hline
$A$矩阵访存 & 合并性好 & 合并性好 & 合并性好 \\
\hline
$B$矩阵访存 & 较好 & 非合并，性能差 & 合并性好 \\
\hline
实现复杂度 & 低 & 低 & 高 \\
\hline
共享内存效率 & 较高 & 较低 & 最高 \\
\hline
相对性能（N=2048） & 1$\times$（基线） & 0.3--0.5$\times$ & 2--3$\times$ \\
\hline
\end{tabularx}
\end{table}

\begin{itemize}
    \item \textbf{按行划分}：实现相对简单（1D grid），对$A$矩阵访问合并性好（每行连续存储），但当$M$较小时并行度不足。适合$M \gg N$的瘦高矩阵乘法场景。
    \item \textbf{按列划分}：同样实现简单（1D grid），但由于矩阵在内存中按行优先存储，访问$B$矩阵的一列意味着跨步读取（stride = $N$），导致非合并访存，严重降低性能。除非配合矩阵转置等预处理，否则不建议使用。
    \item \textbf{按数据块划分}：实现略微复杂（需要2D grid和更精细的索引计算），但提供了最佳的并行度（$O(N^2/\text{tile\_size}^2)$个线程块）、最好的访存合并性（两个矩阵的tile访问都连续），以及最高的共享内存利用率。是通用矩阵乘法的标准实现方式，也是cuBLAS等库的底层策略。
\end{itemize}

    \item \textbf{何时应该优先考虑使用哪种存储？}\par
    回答：

GPU存储层次的选择应遵循"尽可能使用更快、更靠近计算单元的存储"原则：

\begin{enumerate}
    \item \textbf{全局内存（Global Memory）}：仅在数据首次加载和最终写回时使用。应尽量减少运行时的全局内存访问。适合存储整个矩阵，但计算中应避免频繁读取。

    \item \textbf{共享内存（Shared Memory）}：当同一个block内的多个线程需要重复使用相同数据时优先使用。在矩阵乘法中，$A$的每一行和$B$的每一列都被多个线程共享，因此使用共享内存缓存tile极为有效。使用共享内存的典型场景：\textbf{(1)} 数据存在线程间复用；\textbf{(2)} 访问模式需要重排以改善全局内存合并性；\textbf{(3)} 需要线程间通信（如规约操作）。

    \item \textbf{寄存器（Register）}：当单个线程需要反复使用少量数据时优先使用。寄存器数量有限（V100每线程最多255个），应分配给最频繁访问的变量。在矩阵乘法中，累加器\texttt{c}和从共享内存预取的\texttt{a\_reg}、\texttt{b\_reg}天然适合放在寄存器中。寄存器分块通过增加每线程的寄存器用量来换取更高的计算访存比，在算术强度足够高时效果显著。

    \item \textbf{选择策略}：从外到内，数据应逐步从全局内存迁移到共享内存再到寄存器。每一步迁移都应带来数据重用或访存模式改善。对于矩阵乘法：\textbf{先使用共享内存}消除冗余的全局访存；\textbf{再使用寄存器}进一步提高计算密度。但需注意——寄存器是稀缺资源，过度使用会降低占用率，需要在计算密度和并行度之间权衡。
\end{enumerate}

\end{itemize}

\end{document}
